% Options for packages loaded elsewhere
\PassOptionsToPackage{unicode}{hyperref}
\PassOptionsToPackage{hyphens}{url}
\PassOptionsToPackage{dvipsnames,svgnames,x11names}{xcolor}
\documentclass[
  11pt,
  a4paper,
]{ctexart}
\usepackage{xcolor}
\usepackage[top=2.2cm,bottom=2.2cm,left=2.5cm,right=2.5cm]{geometry}
\usepackage{amsmath,amssymb}
\setcounter{secnumdepth}{-\maxdimen} % remove section numbering
\usepackage{iftex}
\ifPDFTeX
  \usepackage[T1]{fontenc}
  \usepackage[utf8]{inputenc}
  \usepackage{textcomp} % provide euro and other symbols
\else % if luatex or xetex
  \usepackage{unicode-math} % this also loads fontspec
  \defaultfontfeatures{Scale=MatchLowercase}
  \defaultfontfeatures[\rmfamily]{Ligatures=TeX,Scale=1}
\fi
\usepackage{lmodern}
\ifPDFTeX\else
  % xetex/luatex font selection
\fi
% Use upquote if available, for straight quotes in verbatim environments
\IfFileExists{upquote.sty}{\usepackage{upquote}}{}
\IfFileExists{microtype.sty}{% use microtype if available
  \usepackage[]{microtype}
  \UseMicrotypeSet[protrusion]{basicmath} % disable protrusion for tt fonts
}{}
\makeatletter
\@ifundefined{KOMAClassName}{% if non-KOMA class
  \IfFileExists{parskip.sty}{%
    \usepackage{parskip}
  }{% else
    \setlength{\parindent}{0pt}
    \setlength{\parskip}{6pt plus 2pt minus 1pt}}
}{% if KOMA class
  \KOMAoptions{parskip=half}}
\makeatother
\setlength{\emergencystretch}{3em} % prevent overfull lines
\providecommand{\tightlist}{%
  \setlength{\itemsep}{0pt}\setlength{\parskip}{0pt}}
\usepackage{amsmath,amssymb,mathtools}
\usepackage{xcolor}
\usepackage{fancyhdr}
\usepackage{enumitem}
\usepackage{booktabs,longtable,array}
\usepackage{microtype}
\definecolor{BookNavy}{HTML}{173F4F}
\definecolor{BookTeal}{HTML}{2F7664}
\definecolor{BookMuted}{HTML}{687B78}
\setlength{\parindent}{2em}
\setlength{\parskip}{0.35em}
\linespread{1.25}
\setlist{itemsep=0.25em,topsep=0.4em,leftmargin=2.2em}
\pagestyle{fancy}
\fancyhf{}
\fancyhead[R]{\small\color{BookMuted}多臂老虎机与重尾在线学习 · 数学学习笔记}
\fancyfoot[C]{\small\color{BookMuted}--\ \thepage\ --}
\renewcommand{\headrulewidth}{0pt}
\renewcommand{\footrulewidth}{0pt}
\setlength{\headheight}{15pt}
\ctexset{
  section={format=\Large\bfseries\color{BookNavy},beforeskip=2.2ex,afterskip=1.2ex},
  subsection={format=\large\bfseries\color{BookTeal},beforeskip=1.8ex,afterskip=0.8ex},
  subsubsection={format=\normalsize\bfseries\color{BookNavy},beforeskip=1.3ex,afterskip=0.5ex}
}
\usepackage{bookmark}
\IfFileExists{xurl.sty}{\usepackage{xurl}}{} % add URL line breaks if available
\urlstyle{same}
\hypersetup{
  pdftitle={多臂老虎机与重尾在线学习},
  colorlinks=true,
  linkcolor={BookTeal},
  filecolor={Maroon},
  citecolor={Blue},
  urlcolor={BookTeal},
  pdfcreator={LaTeX via pandoc}}

\title{多臂老虎机与重尾在线学习}
\usepackage{etoolbox}
\makeatletter
\providecommand{\subtitle}[1]{% add subtitle to \maketitle
  \apptocmd{\@title}{\par {\large #1 \par}}{}{}
}
\makeatother
\subtitle{从概率不等式到 Robust UCB 与 uniINF}
\author{}
\date{zyx学习笔记}

\begin{document}
\maketitle

\thispagestyle{empty}
\vfill
\begin{center}
{\color{BookMuted}\large 从定义、定理与证明出发的在线学习入门}
\end{center}
\vfill
\clearpage
\setcounter{tocdepth}{2}
\tableofcontents
\clearpage



\begin{center}\rule{0.5\linewidth}{0.5pt}\end{center}

\subsection{阅读说明}\label{ux9605ux8bfbux8bf4ux660e}

这份笔记讨论一个贯穿始终的问题：当我们必须一边行动、一边学习未知环境时，怎样在``探索未知选项''和``利用当前最佳选项''之间取得平衡？

全文的逻辑路线是

\[
\begin{aligned}
\text{概率与集中不等式}
&\longrightarrow \text{随机多臂老虎机}
\longrightarrow \text{ETC与UCB}\\
&\longrightarrow \text{对抗型老虎机与EXP3}
\longrightarrow \text{重尾反馈与稳健估计}\\
&\longrightarrow \text{Robust UCB}
\longrightarrow \text{uniINF}\\
&\longrightarrow \text{持续更新中}.
\end{aligned}
\]


\begin{center}\rule{0.5\linewidth}{0.5pt}\end{center}

\section{第一章
概率论准备：从尾概率到置信区间}\label{ux7b2cux4e00ux7ae0-ux6982ux7387ux8bbaux51c6ux5907ux4eceux5c3eux6982ux7387ux5230ux7f6eux4fe1ux533aux95f4}

\subsection{1.1
随机变量、期望和方差}\label{ux968fux673aux53d8ux91cfux671fux671bux548cux65b9ux5dee}

\subsubsection{定义 1.1（期望）}\label{ux5b9aux4e49-1.1ux671fux671b}

设离散随机变量 \(X\) 取值为 \(x_1,x_2,\ldots\)，则

\[
\mathbb E[X]=\sum_jx_j\mathbb P(X=x_j),
\]

只要右侧绝对收敛。连续型随机变量的期望定义为相应积分。

期望描述大量重复实验中观测值的长期平均水平。

\subsubsection{定义 1.2（方差）}\label{ux5b9aux4e49-1.2ux65b9ux5dee}

设 \(\mu=\mathbb E[X]\)，则

\[
\operatorname{Var}(X)=\mathbb E[(X-\mu)^2].
\]

若二阶矩有限，则

\[
\operatorname{Var}(X)=\mathbb E[X^2]-\mu^2.
\]

方差描述随机变量围绕均值的平均平方波动。

\subsubsection{定义
1.3（示性函数）}\label{ux5b9aux4e49-1.3ux793aux6027ux51fdux6570}

对事件 \(A\)，定义

\[
\mathbf 1_A(\omega)=
\begin{cases}
1,&\omega\in A,\\
0,&\omega\notin A.
\end{cases}
\]

因此

\[
\mathbb E[\mathbf1_A]=\mathbb P(A).
\]

示性函数把``事件是否发生''转化为数值变量，是概率证明中非常常见的桥梁。

\subsection{1.2 Markov 不等式}\label{markov-ux4e0dux7b49ux5f0f}

\subsubsection{定理 1.1（Markov
不等式）}\label{ux5b9aux7406-1.1markov-ux4e0dux7b49ux5f0f}

若 \(X\ge 0\) 几乎处处成立，且 \(a>0\)，则

\[
\mathbb P(X\ge a)\le \frac{\mathbb E[X]}{a}.
\]

\subsubsection{证明}\label{ux8bc1ux660e}

逐点有

\[
X\ge a\mathbf1_{\{X\ge a\}}.
\]

两边取期望，得到

\[
\mathbb E[X]
\ge a\mathbb E[\mathbf1_{\{X\ge a\}}]
=a\mathbb P(X\ge a).
\]

两边除以 \(a\)，即得结论。\(\square\)

\subsubsection{直观解释}\label{ux76f4ux89c2ux89e3ux91ca}

如果一个非负随机变量的平均值很小，它就不可能以很高概率取到特别大的值。Markov
不等式只使用期望，因此适用范围广，但通常较松。

\subsection{1.3 Chebyshev 不等式}\label{chebyshev-ux4e0dux7b49ux5f0f}

\subsubsection{定理 1.2（Chebyshev
不等式）}\label{ux5b9aux7406-1.2chebyshev-ux4e0dux7b49ux5f0f}

设 \(X\) 的均值为 \(\mu\)，方差为 \(\sigma^2<\infty\)。对任意 \(a>0\)，

\[
\mathbb P(|X-\mu|\ge a)
\le \frac{\sigma^2}{a^2}.
\]

\subsubsection{证明}\label{ux8bc1ux660e-1}

令非负随机变量

\[
Y=(X-\mu)^2.
\]

注意到

\[
\{|X-\mu|\ge a\}=\{Y\ge a^2\}.
\]

对 \(Y\) 使用 Markov 不等式：

\[
\mathbb P(|X-\mu|\ge a)
=\mathbb P(Y\ge a^2)
\le\frac{\mathbb E[Y]}{a^2}
=\frac{\sigma^2}{a^2}.
\]

证毕。\(\square\)

\subsubsection{推论 1.1（样本均值的 Chebyshev
界）}\label{ux63a8ux8bba-1.1ux6837ux672cux5747ux503cux7684-chebyshev-ux754c}

设 \(X_1,\ldots,X_n\) 独立同分布，均值为 \(\mu\)，方差为
\(\sigma^2\)。令

\[
\bar X_n=\frac1n\sum_{t=1}^nX_t.
\]

则

\[
\mathbb P(|\bar X_n-\mu|\ge a)
\le\frac{\sigma^2}{na^2}.
\]

\subsubsection{证明}\label{ux8bc1ux660e-2}

由独立性，

\[
\operatorname{Var}(\bar X_n)
=\frac1{n^2}\sum_{t=1}^n\operatorname{Var}(X_t)
=\frac{\sigma^2}{n}.
\]

再对 \(\bar X_n\) 应用 Chebyshev 不等式即可。\(\square\)

\subsection{1.4 Sub-Gaussian
随机变量}\label{sub-gaussian-ux968fux673aux53d8ux91cf}

\subsubsection{定义
1.4（Sub-Gaussian）}\label{ux5b9aux4e49-1.4sub-gaussian}

若存在 \(v>0\)，使得对任意 \(\lambda\in\mathbb R\)，

\[
\mathbb E\exp\{\lambda(X-\mathbb E X)\}
\le \exp\left(\frac{v\lambda^2}{2}\right),
\]

则称 \(X\) 是方差因子为 \(v\) 的 sub-Gaussian 随机变量。

\subsubsection{定理 1.3（Sub-Gaussian
尾界）}\label{ux5b9aux7406-1.3sub-gaussian-ux5c3eux754c}

若 \(X\) 满足上述条件，则对任意 \(a>0\)，

\[
\mathbb P(X-\mathbb EX\ge a)
\le\exp\left(-\frac{a^2}{2v}\right).
\]

\subsubsection{证明}\label{ux8bc1ux660e-3}

对任意 \(\lambda>0\)，由 Markov 不等式，

\[
\begin{aligned}
\mathbb P(X-\mathbb EX\ge a)
&=\mathbb P\big(e^{\lambda(X-\mathbb EX)}\ge e^{\lambda a}\big)\\
&\le e^{-\lambda a}\mathbb E e^{\lambda(X-\mathbb EX)}\\
&\le\exp\left(-\lambda a+\frac{v\lambda^2}{2}\right).
\end{aligned}
\]

右端关于 \(\lambda\) 的最小值在 \(\lambda=a/v\) 处取得。代入即得

\[
\mathbb P(X-\mathbb EX\ge a)
\le e^{-a^2/(2v)}.
\]

证毕。\(\square\)

\subsubsection{有限方差为什么弱于
sub-Gaussian？}\label{ux6709ux9650ux65b9ux5deeux4e3aux4ec0ux4e48ux5f31ux4e8e-sub-gaussian}

Sub-Gaussian
条件给出指数衰减的尾概率，并能推出所有阶矩有限。仅有有限方差只要求

\[
\mathbb E[(X-\mu)^2]<\infty,
\]

仍允许尾概率按照多项式速度衰减。因此

\[
\text{sub-Gaussian}\Longrightarrow\text{有限方差},
\]

反向一般不成立。

\subsection{1.5 置信区间}\label{ux7f6eux4fe1ux533aux95f4}

若估计量 \(\widehat\mu_n\) 满足

\[
\mathbb P\big(|\widehat\mu_n-\mu|\le r(n,\delta)\big)
\ge1-\delta,
\]

则区间

\[
[\widehat\mu_n-r(n,\delta),\widehat\mu_n+r(n,\delta)]
\]

称为置信水平至少为 \(1-\delta\) 的置信区间，\(r(n,\delta)\)
称为置信半径。

置信半径应具有两个方向：样本数 \(n\)
越大，半径越小；要求的置信水平越高，即 \(\delta\) 越小，半径越大。

\subsection{本章关键疑问与解答}\label{ux672cux7ae0ux5173ux952eux7591ux95ee}

\textbf{疑问一：为什么 Markov、Chebyshev 和 sub-Gaussian 尾界不是可以互相替换的三个公式？}

\textbf{解答：}它们使用的信息强弱不同。Markov 只使用非负性和一阶矩；Chebyshev 使用有限方差；sub-Gaussian 条件控制矩母函数，从而给出指数尾界。假设越强，尾概率通常下降得越快。只有有限方差时，不能未经证明就套用指数尾界，这也是后文重尾估计必须更换工具的根本原因。

\textbf{疑问二：sub-Gaussian 尾界的证明为什么要引入参数 \(\lambda\)，最后还要对它优化？}

\textbf{解答：}指数变换把事件 \(X-\mathbb EX\ge a\) 转化为一个可用 Markov 不等式处理的非负随机变量事件。每个 \(\lambda>0\) 都产生一个合法上界，而 \(\lambda\) 只是证明者可以选择的参数；取最小值是在所有合法上界中选最紧者。这里的优化不是算法步骤，而是证明步骤。

\textbf{疑问三：固定样本量的置信区间为什么不能直接保证一个自适应算法在所有时刻都正确？}

\textbf{解答：}老虎机算法会根据历史决定何时、选择哪个臂，因此样本量本身是随机的，而且我们关心很多时刻和很多臂同时不失败。证明时通常要先对固定样本数建立界，再对样本数、时间和臂取并集界，或直接使用适用于停止时刻的置信序列。忽略这一层会把“某一次高概率正确”误写成“整个运行过程高概率正确”。

\begin{center}\rule{0.5\linewidth}{0.5pt}\end{center}

\section{第二章
随机平稳多臂老虎机}\label{ux7b2cux4e8cux7ae0-ux968fux673aux5e73ux7a33ux591aux81c2ux8001ux864eux673a}

\subsection{2.1 问题定义}\label{ux95eeux9898ux5b9aux4e49}

\subsubsection{定义
2.1（随机平稳多臂老虎机）}\label{ux5b9aux4e49-2.1ux968fux673aux5e73ux7a33ux591aux81c2ux8001ux864eux673a}

设有 \(K\) 个臂。臂 \(i\) 对应未知奖励分布 \(\nu_i\)，均值为

\[
\mu_i=\mathbb E_{X\sim\nu_i}[X].
\]

第 \(t\) 轮，学习者选择臂
\(A_t\in\{1,\ldots,K\}\)，随后只观察该臂产生的奖励
\(X_{t,A_t}\)。对同一臂的各次奖励独立同分布，且分布不随时间变化。

``随机''表示奖励具有随机性；``平稳''表示每个臂的奖励分布随时间保持不变。

\subsubsection{定义 2.2（最优臂与
gap）}\label{ux5b9aux4e49-2.2ux6700ux4f18ux81c2ux4e0e-gap}

令

\[
\mu^*=\max_{1\le i\le K}\mu_i.
\]

任取满足 \(\mu_{i^*}=\mu^*\) 的臂称为最优臂。定义臂 \(i\) 的次优差距

\[
\Delta_i=\mu^*-\mu_i\ge0.
\]

\subsubsection{定义
2.3（伪遗憾）}\label{ux5b9aux4e49-2.3ux4f2aux9057ux61be}

前 \(T\) 轮的期望伪遗憾定义为

\[
R_T=T\mu^*-\mathbb E\left[\sum_{t=1}^T\mu_{A_t}\right].
\]

令

\[
T_i(T)=\sum_{t=1}^T\mathbf1_{\{A_t=i\}}
\]

表示臂 \(i\) 被选择的次数，则

\[
R_T=\sum_{i=1}^K\Delta_i\mathbb E[T_i(T)].
\]

\subsubsection{证明}\label{ux8bc1ux660e-4}

因为每轮选择臂 \(A_t\) 所造成的期望损失是

\[
\mu^*-\mu_{A_t}=\sum_{i=1}^K\Delta_i\mathbf1_{\{A_t=i\}},
\]

所以对时间求和并取期望：

\[
\begin{aligned}
R_T
&=\mathbb E\sum_{t=1}^T(\mu^*-\mu_{A_t})\\
&=\sum_{i=1}^K\Delta_i\mathbb E\sum_{t=1}^T\mathbf1_{\{A_t=i\}}\\
&=\sum_{i=1}^K\Delta_i\mathbb E[T_i(T)].
\end{aligned}
\]

证毕。\(\square\)

\subsection{2.2 探索与利用}\label{ux63a2ux7d22ux4e0eux5229ux7528}

若只选择当前经验均值最高的臂，早期噪声可能导致永久误判；若始终随机探索，又无法利用已经获得的信息。老虎机问题的核心就是在两者之间取得平衡。

\subsubsection{Random 策略}\label{random-ux7b56ux7565}

每轮均匀随机选择一个臂。它保证探索，却通常产生线性遗憾。

\subsubsection{Greedy 策略}\label{greedy-ux7b56ux7565}

每轮选择当前经验均值最大的臂。它可能被早期偶然高奖励锁定。

\subsubsection{\texorpdfstring{\(\varepsilon\)-greedy
策略}{\textbackslash varepsilon-greedy 策略}}\label{varepsilon-greedy-ux7b56ux7565}

以概率 \(1-\varepsilon\) 选择经验最优臂，以概率 \(\varepsilon\)
均匀探索。它比 Greedy 稳健，但固定 \(\varepsilon\)
会造成持续探索，通常仍产生与 \(T\) 成正比的探索损失。

\subsection{2.3 Bernoulli 环境}\label{bernoulli-ux73afux5883}

若臂 \(i\) 的奖励满足

\[
X_{t,i}\sim\operatorname{Bernoulli}(p_i),
\]

则奖励只能为 \(0\) 或 \(1\)，且

\[
\mu_i=\mathbb E[X_{t,i}]=p_i.
\]

因此学习最佳臂等价于识别成功概率最大的臂。

\subsection{本章关键疑问与解答}\label{ux672cux7ae0ux5173ux952eux7591ux95ee-1}

\textbf{疑问一：为什么遗憾可以写成 \(R_T=\sum_i\Delta_i\mathbb E[T_i(T)]\)，这个式子真正告诉了我们什么？}

\textbf{解答：}每次选择臂 \(i\) 都付出期望差距 \(\Delta_i\)，因此总遗憾等于“每次选错的代价”乘以“选错的次数”再求和。算法分析由此转化为控制次优臂的选择次数。它也说明小 gap 的臂虽然单次损失小，却通常更难识别，可能被选择更多次。

\textbf{疑问二：为什么固定 \(\varepsilon>0\) 的 \(\varepsilon\)-greedy 通常不能得到次线性遗憾？}

\textbf{解答：}即使算法早已确定最优臂，每轮仍以固定概率探索，并以正概率选择次优臂。因此到第 \(T\) 轮，非必要探索次数与 \(T\) 同阶，累计遗憾也与 \(T\) 同阶。要获得次线性遗憾，探索强度必须随信息增加而衰减，或由置信度自适应决定。

\textbf{疑问三：随机平稳假设在分析中承担了什么作用？}

\textbf{解答：}它保证同一臂过去的样本能够估计未来的同一个均值。若分布随时间漂移，旧数据可能不再代表当前环境；若奖励依赖学习者过去的行动，独立同分布的集中论证也会失效。因此平稳随机老虎机不是“奖励每次固定”，而是“产生奖励的统计规律固定”。

\begin{center}\rule{0.5\linewidth}{0.5pt}\end{center}

\section{第三章
Explore-Then-Commit}\label{ux7b2cux4e09ux7ae0-explore-then-commit}

\subsection{3.1 算法定义}\label{ux7b97ux6cd5ux5b9aux4e49}

\subsubsection{定义 3.1（ETC）}\label{ux5b9aux4e49-3.1etc}

预先选择探索长度 \(m\)。先将每个臂各选择 \(m\) 次，计算经验均值

\[
\widehat\mu_i(m)=\frac1m\sum_{s=1}^mX_{i,s}.
\]

随后选择经验均值最大的臂，并在剩余轮次中始终选择该臂。

\subsection{3.2
为什么探索长度需要平衡？}\label{ux4e3aux4ec0ux4e48ux63a2ux7d22ux957fux5ea6ux9700ux8981ux5e73ux8861}

\begin{itemize}
\tightlist
\item
  \(m\) 太小：识别错误的概率较大，之后可能长期选择错误臂；
\item
  \(m\) 太大：即使已经足以识别最优臂，仍被迫继续探索次优臂。
\end{itemize}

考虑两臂且奖励为 sub-Gaussian。若 gap 为
\(\Delta\)，经验均值判断错误要求至少一个臂发生约 \(\Delta/2\)
的偏差。集中不等式给出

\[
\mathbb P(\text{识别错误})
\lesssim e^{-cm\Delta^2}.
\]

因此遗憾可粗略分解为

\[
R_T
\lesssim
m\Delta+T\Delta e^{-cm\Delta^2}.
\]

第一项是强制探索损失，第二项是识别错误后的长期损失。令第二项不至于过大，通常需要

\[
m\asymp\frac{\log T}{\Delta^2}.
\]

ETC展示了探索与利用的基本权衡，但它需要预先确定探索长度，且通常需要了解
gap 的尺度。

\iffalse

\begin{enumerate}
\def\labelenumi{\arabic{enumi}.}
\tightlist
\item
  在上述粗略遗憾分解中，说明为什么增大 \(m\)
  会使第一项增大、第二项减小。
\item
  若 \(\Delta\) 很小，为什么 ETC 所需探索次数会迅速增加？
\end{enumerate}

\subsection{本章习题解答}\label{ux672cux7ae0ux4e60ux9898ux89e3ux7b54-2}

\begin{enumerate}
\def\labelenumi{\arabic{enumi}.}
\tightlist
\item
  第一项 \(m\Delta\) 随 \(m\) 线性增大；第二项含 \(e^{-cm\Delta^2}\)，随
  \(m\) 指数减小。
\item
  两臂均值越接近越难区分，而 \(m\asymp(\log T)/\Delta^2\)，故
  \(\Delta\downarrow0\) 时所需样本数按 \(1/\Delta^2\) 境长。
\end{enumerate}

\begin{center}\rule{0.5\linewidth}{0.5pt}\end{center}

\fi
\subsection{本章关键疑问与解答}\label{ux672cux7ae0ux5173ux952eux7591ux95ee-2}

\textbf{疑问一：ETC 的探索长度为什么必须同时控制探索损失和误判损失？}

\textbf{解答：}增加探索长度会线性增加强制尝试次优臂的代价，却会指数降低均值排序错误的概率。过早停止可能导致剩余所有轮次都选择错误臂；过晚停止则把已经足够清楚的问题继续当作未知问题处理。最优尺度来自两种代价的平衡，而不是单独让某一项最小。

\textbf{疑问二：为什么不知道 gap 会成为 ETC 的结构性困难？}

\textbf{解答：}区分两臂所需样本量约为 \(\Delta^{-2}\log T\)，而 ETC 必须在看完数据前预先决定探索长度。设得适合大 gap，会在小 gap 时过早承诺；设得适合极小 gap，又会在容易问题上浪费样本。这促使我们转向能根据置信半径动态继续或停止探索的算法。

\textbf{疑问三：ETC 为什么比 UCB 更容易受到一次错误决定的长期影响？}

\textbf{解答：}ETC 的 commit 阶段不再修正选择，所以探索结束时的排序错误会持续到终点。UCB 每轮重新计算指标，新数据能够逐步纠正早期误判。两者的差异不是是否探索，而是探索与利用能否在整个过程中持续交替。

\begin{center}\rule{0.5\linewidth}{0.5pt}\end{center}

\section{第四章
UCB：乐观面对不确定性}\label{ux7b2cux56dbux7ae0-ucbux4e50ux89c2ux9762ux5bf9ux4e0dux786eux5b9aux6027}

\subsection{4.1 算法思想}\label{ux7b97ux6cd5ux601dux60f3}

UCB
不把探索和利用分成两个固定阶段，而是在每一轮同时考虑当前估计和不确定性。

\subsubsection{定义 4.1（UCB
指标）}\label{ux5b9aux4e49-4.1ucb-ux6307ux6807}

设在第 \(t\) 轮前臂 \(i\) 已被选择 \(N_i(t-1)\) 次，经验均值为
\(\widehat\mu_i(t-1)\)。定义

\[
B_i(t)=\widehat\mu_i(t-1)+r_i(t),
\]

其中 \(r_i(t)\) 是置信半径。第 \(t\) 轮选择

\[
A_t\in\arg\max_iB_i(t).
\]

指标的左项负责利用，右项负责探索。样本少时半径大；样本增多后半径缩小。

对于未选择过的臂，可规定
\(B_i(t)=+\infty\)，从而保证每个臂至少被尝试一次。

\subsection{4.2
次优臂被选择的三种原因}\label{ux6b21ux4f18ux81c2ux88abux9009ux62e9ux7684ux4e09ux79cdux539fux56e0}

令 \(i^*\) 为最优臂，\(i\) 为次优臂。若算法选择了
\(i\)，则至少发生下列事件之一：

\begin{enumerate}
\def\labelenumi{\arabic{enumi}.}
\tightlist
\item
  最优臂被低估：\(B_{i^*}(t)<\mu^*\)；
\item
  次优臂被高估：\(\widehat\mu_i(t)>\mu_i+r_i(t)\)；
\item
  次优臂样本不足：\(2r_i(t)>\Delta_i\)。
\end{enumerate}

\subsubsection{证明}\label{ux8bc1ux660e-5}

反设三个事件都不发生。则

\[
B_{i^*}(t)\ge\mu^*.
\]

又因为第二、三个事件不发生，

\[
B_i(t)=\widehat\mu_i(t)+r_i(t)
\le\mu_i+2r_i(t)
\le\mu_i+\Delta_i
=\mu^*.
\]

因此 \(B_i(t)\le B_{i^*}(t)\)，次优臂不可能严格优先于最优臂。\(\square\)

\subsection{4.3
经典遗憾界的证明骨架}\label{ux7ecfux5178ux9057ux61beux754cux7684ux8bc1ux660eux9aa8ux67b6}

\subsubsection*{补充证明：sub-Gaussian 置信半径从何而来？}

这不是凭空规定的公式，而是从样本均值的偏差概率上界反解出来的。以下先处理固定样本数，再处理老虎机中随机的拉取次数。全文的 \(\log\) 均为自然对数。

\paragraph{一、假设与定义}
固定臂 \(i\)，将该臂依次被拉取时的奖励记为
\(X_{i,1},X_{i,2},\ldots\)。假设它们独立同分布、均值为 \(\mu_i\)，且满足
\[
\mathbb E\!\left[e^{\lambda(X_{i,k}-\mu_i)}\right]
\le e^{\sigma^2\lambda^2/2},
\qquad \lambda\in\mathbb R,\quad \sigma^2>0.
\]
这称为 \(\sigma^2\)-sub-Gaussian 条件；\(\sigma^2\) 是方差代理，不必等于真实方差。为简洁起见，假设各臂共用此上界。前 \(n\) 个奖励的样本均值为
\[
\widehat\mu_{i,n}=\frac1n\sum_{k=1}^n X_{i,k}.
\]
目标是找出半径 \(r\)，使估计偏离真实均值超过 \(r\) 的概率足够小。

\paragraph{二、定理：固定样本数的集中不等式}
对任意确定的 \(n\ge1\) 和 \(r>0\)，有
\[
\boxed{\mathbb P\!\left(
|\widehat\mu_{i,n}-\mu_i|\ge r
\right)\le 2\exp\!\left(-\frac{nr^2}{2\sigma^2}\right).}
\]

证明第一步：利用独立性控制样本均值的矩母函数：
\[
\begin{aligned}
\mathbb E e^{\lambda(\widehat\mu_{i,n}-\mu_i)}
&=\prod_{k=1}^n
\mathbb E e^{(\lambda/n)(X_{i,k}-\mu_i)}\\
&\le\prod_{k=1}^n
\exp\!\left(\frac{\sigma^2\lambda^2}{2n^2}\right)
=\exp\!\left(\frac{\sigma^2\lambda^2}{2n}\right).
\end{aligned}
\]
所以取平均后，方差代理从 \(\sigma^2\) 缩小为 \(\sigma^2/n\)。

第二步：对任意 \(\lambda>0\)，对非负随机变量
\(e^{\lambda(\widehat\mu_{i,n}-\mu_i)}\) 使用 Markov 不等式：
\[
\begin{aligned}
\mathbb P(\widehat\mu_{i,n}-\mu_i\ge r)
&=\mathbb P\!\left(e^{\lambda(\widehat\mu_{i,n}-\mu_i)}
                    \ge e^{\lambda r}\right)\\
&\le e^{-\lambda r}
       \mathbb E e^{\lambda(\widehat\mu_{i,n}-\mu_i)}\\
&\le \exp\!\left(-\lambda r+
                      \frac{\sigma^2\lambda^2}{2n}\right).
\end{aligned}
\]

第三步：选择使上界最小的 \(\lambda\)。将指数配方：
\[
-\lambda r+\frac{\sigma^2\lambda^2}{2n}
=\frac{\sigma^2}{2n}
  \left(\lambda-\frac{nr}{\sigma^2}\right)^2
 -\frac{nr^2}{2\sigma^2}.
\]
取 \(\lambda=nr/\sigma^2>0\)，得到上尾界
\[
\mathbb P(\widehat\mu_{i,n}-\mu_i\ge r)
\le e^{-nr^2/(2\sigma^2)}.
\]
对负偏差重复同样的证明，得到下尾界。两侧事件取并集，概率上界相加，即得定理。证毕。

\paragraph{三、反解置信半径}
给定允许的失败概率 \(\delta\in(0,1)\)，只需要求
\[
2e^{-nr^2/(2\sigma^2)}\le\delta.
\]
取对数、移项并取正平方根，可取
\[
\boxed{r(n,\delta)=
\sqrt{\frac{2\sigma^2\log(2/\delta)}{n}}.}
\]
于是 \([\widehat\mu_{i,n}-r,\widehat\mu_{i,n}+r]\)
以至少 \(1-\delta\) 的概率覆盖 \(\mu_i\)。
样本数越大，半径按 \(1/\sqrt n\) 缩小；要求的失败概率越小，半径越大。

对于 \(t\ge2\)，若令固定样本数处的失败概率为
\(\delta_t=2t^{-\beta}\)，其中 \(\beta>1\)，便有
\[
r(n,t)=\sqrt{\frac{2\sigma^2\beta\log t}{n}}
       =\sqrt{\frac{c\log t}{n}},
\qquad c=2\sigma^2\beta.
\]
因此 \(\log t\) 来自让错误概率随轮数按多项式速度下降，而不是一个额外的经验假设。

\paragraph{四、关键细节：拉取次数是随机变量}
在第 \(t\) 轮决策前，
\[
\widehat\mu_i(t-1)=
\widehat\mu_{i,N_i(t-1)}.
\]
算法根据过去奖励选择动作，故 \(N_i(t-1)\) 与奖励相关。
不能直接在条件 \(N_i(t-1)=n\) 下照搬固定样本数的独立同分布集中界：
这个条件本身可能筛选了奖励轨迹。

解决方法是同时控制所有可能的样本数。固定 \(t\ge2\)，定义
\[
E_{i,t}=
\left\{N_i(t-1)\ge1,\
|\widehat\mu_{i,N_i(t-1)}-\mu_i|
>\sqrt{\frac{c\log t}{N_i(t-1)}}\right\}.
\]
随机次数处若出错，则某个确定样本数 \(1\le n\le t-1\) 处必然出错，因此
\[
E_{i,t}\subseteq
\bigcup_{n=1}^{t-1}
\left\{|\widehat\mu_{i,n}-\mu_i|
>\sqrt{\frac{c\log t}{n}}\right\}.
\]
并集界不要求这些事件相互独立。由固定样本数的定理，
\[
\begin{aligned}
\mathbb P(E_{i,t})
&\le\sum_{n=1}^{t-1}
2\exp\!\left[-\frac{n}{2\sigma^2}
                  \frac{c\log t}{n}\right]\\
&=2(t-1)t^{-c/(2\sigma^2)}
\le 2t^{\,1-c/(2\sigma^2)}.
\end{aligned}
\]
这里的额外因子 \(t\) 来自对可能样本数取并集，不能漏掉。

例如取 \(c=6\sigma^2\)（即 \(\beta=3\)），就有
\[
\mathbb P(E_{i,t})\le\frac2{t^2},
\qquad
\sum_{t=2}^{\infty}\mathbb P(E_{i,t})<\infty.
\]
若还要求 \(K\) 个臂在同一轮都正确，再对臂取并集，失效概率至多为 \(2K/t^2\)。
这些是每轮的概率界及其可求和性，不是说所有轮次必然都正确。
常数 \(6\sigma^2\) 是便于证明的充分选择，并非最优常数。

\paragraph{五、回到 UCB 与边界情形}
由此得到本节使用的形式：
\[
\boxed{r_i(t)=\sqrt{\frac{c\log t}{N_i(t-1)}}.}
\]
若 \(N_i(t-1)=0\)，不能做除法：通常先将每个臂拉取一次，或者把未拉取臂的 UCB 指数设为 \(+\infty\)。
在上述简单并集证明中，取 \(c=6\sigma^2\) 可直接接上后面的“失败概率按时间可求和”。

概括这条链：
sub-Gaussian 条件给出指数型偏差界；反解失败概率得到置信半径；
选择随时间下降的失败概率产生 \(\log t\)；
再用并集界处理自适应拉取造成的随机样本数。

\paragraph{六、自检题与解答}
题目：固定 \(t\) 与 \(c\)，若某个臂的拉取次数从 \(n\) 增加至 \(4n\)，置信半径变成多少？
\[
r(4n,t)=\sqrt{\frac{c\log t}{4n}}
       =\frac12r(n,t).
\]
答案：变为原来的一半。因此若想把半径减半，需要将样本数增加至原来的四倍。

\subsubsection*{继续：由置信半径得到拉取次数界}

若 sub-Gaussian 置信半径为

\[
r_i(t)=\sqrt{\frac{c\log t}{N_i(t-1)}},
\]

则要使 \(2r_i(t)\le\Delta_i\)，只需

\[
N_i(t-1)\ge\frac{4c\log t}{\Delta_i^2}.
\]

置信区间失败事件可通过令其概率按 \(t^{-2}\)
或更快速度下降而求和控制。因此

\[
\mathbb E[T_i(T)]
=O\left(\frac{\log T}{\Delta_i^2}\right).
\]

再利用

\[
R_T=\sum_i\Delta_i\mathbb E[T_i(T)],
\]

得到

\[
R_T=O\left(\sum_{i:\Delta_i>0}\frac{\log T}{\Delta_i}\right).
\]

\iffalse

\begin{enumerate}
\def\labelenumi{\arabic{enumi}.}
\tightlist
\item
  两个臂的估计均值分别为 \(2.5,2.2\)，置信半径分别为
  \(0.3,0.8\)。计算UCB指标并确定选择哪个臂。
\item
  为什么样本较少的臂即使估计均值稍低，也可能被UCB选择？
\item
  从 \(2\sqrt{c\log T/N_i}\le\Delta_i\) 解出 \(N_i\) 的下界。
\end{enumerate}

\subsection{本章习题解答}\label{ux672cux7ae0ux4e60ux9898ux89e3ux7b54-3}

\begin{enumerate}
\def\labelenumi{\arabic{enumi}.}
\tightlist
\item
  \(B_1=2.8,B_2=3.0\)，因此选择第二个臂。
\item
  因为它的置信半径更大，真实均值仍可能很高；这是``乐观面对不确定性''。
\item
  两边平方并整理，得到 \(N_i\ge4c\log T/\Delta_i^2\)。
\end{enumerate}

\begin{center}\rule{0.5\linewidth}{0.5pt}\end{center}

\fi
\subsection{本章关键疑问与解答}\label{ux672cux7ae0ux5173ux952eux7591ux95ee-3}

\textbf{疑问一：UCB 为什么能同时完成探索和利用，而不需要显式设置探索阶段？}

\textbf{解答：}指标由经验均值和置信半径组成。经验均值高的臂因看起来优秀而被利用；样本少的臂因置信半径大而仍有机会被探索。每选择一次某臂，它的半径就缩小，因此探索会随着证据积累自动减少。

\textbf{疑问二：证明中“次优臂被选择只可能因为三个坏事件之一”为什么是核心？}

\textbf{解答：}当最优臂没有被低估、次优臂没有被高估，且次优臂半径已小于 gap 的一半时，次优臂的指标不可能超过最优臂。于是选择次数被分成两类：采样尚不充分的有限次数，以及置信事件失败的稀有次数。前者给出 \(\log T/\Delta_i^2\) 的样本复杂度，后者通过可求和失败概率控制。

\textbf{疑问三：置信半径中的 \(\log t\) 为什么不是随意添加的惩罚项？}

\textbf{解答：}算法要在许多轮次上反复使用置信结论。把第 \(t\) 轮失败概率设计成约 \(t^{-p}\)（\(p>1\)），才能使全时间失败概率之和有限。反解指数集中界便产生 \(\sqrt{\log t/N_i(t-1)}\)。因此 \(\log t\) 来自“同时控制很多时刻”的概率要求。

\begin{center}\rule{0.5\linewidth}{0.5pt}\end{center}

\section{第五章 对抗型老虎机与
EXP3}\label{ux7b2cux4e94ux7ae0-ux5bf9ux6297ux578bux8001ux864eux673aux4e0e-exp3}

\subsection{5.1
从随机环境到对抗环境}\label{ux4eceux968fux673aux73afux5883ux5230ux5bf9ux6297ux73afux5883}

随机平稳环境假设每个臂的奖励来自固定分布。对抗环境允许每轮的损失向量

\[
\ell_t=(\ell_{t,1},\ldots,\ell_{t,K})
\]

随时间变化，甚至由对手设计。学习者只观察所选臂的损失，这称为 bandit
feedback。

\subsubsection{定义
5.1（对抗型老虎机遗憾）}\label{ux5b9aux4e49-5.1ux5bf9ux6297ux578bux8001ux864eux673aux9057ux61be}

\[
R_T=
\mathbb E\left[\sum_{t=1}^T\ell_{t,A_t}\right]
-\min_{1\le i\le K}\sum_{t=1}^T\ell_{t,i}.
\]

比较对象是事后最好的固定臂，而不是每轮都能变化的动态最优臂。

\subsection{5.2
重要性加权估计}\label{ux91cdux8981ux6027ux52a0ux6743ux4f30ux8ba1}

设第 \(t\) 轮按照概率向量 \(p_t\) 选择臂。定义

\[
\widehat\ell_{t,i}
=\frac{\ell_{t,i}\mathbf1_{\{A_t=i\}}}{p_{t,i}}.
\]

\subsubsection{定理
5.1（无偏性）}\label{ux5b9aux7406-5.1ux65e0ux504fux6027}

在给定历史和损失向量的条件下，

\[
\mathbb E[\widehat\ell_{t,i}]=\ell_{t,i}.
\]

\subsubsection{证明}\label{ux8bc1ux660e-6}

因为 \(\mathbb P(A_t=i)=p_{t,i}\)，所以

\[
\begin{aligned}
\mathbb E[\widehat\ell_{t,i}]
&=p_{t,i}\frac{\ell_{t,i}}{p_{t,i}}
+(1-p_{t,i})\cdot0\\
&=\ell_{t,i}.
\end{aligned}
\]

证毕。\(\square\)

除以 \(p_{t,i}\) 修正了``只有以概率 \(p_{t,i}\)
才能观察到该臂''的抽样偏差，但当 \(p_{t,i}\)
很小时，估计量的波动会很大。

\subsection{5.3 EXP3
的基本结构}\label{exp3-ux7684ux57faux672cux7ed3ux6784}

初始化每个臂权重 \(w_{1,i}=1\)。第 \(t\) 轮：

\begin{enumerate}
\def\labelenumi{\arabic{enumi}.}
\tightlist
\item
  由权重构造概率 \(p_{t,i}\)；
\item
  按 \(p_t\) 随机选择臂 \(A_t\)；
\item
  观察 \(\ell_{t,A_t}\)；
\item
  构造 \(\widehat\ell_{t,i}\)；
\item
  更新
\end{enumerate}

\[
w_{t+1,i}=w_{t,i}\exp(-\eta\widehat\ell_{t,i}).
\]

EXP3不是根据一次奖励断言某个臂永远更好，而是不断调整概率。即使某次正反馈之后下一次出现负反馈，后续权重仍会被重新修正。

在有界损失条件下，合适选择学习率可得到典型遗憾量级

\[
R_T=O(\sqrt{TK\log K}).
\]

\iffalse

\begin{enumerate}
\def\labelenumi{\arabic{enumi}.}
\tightlist
\item
  若 \(p_{t,i}=0.1\)，观察到损失 \(0.5\)，则重要性加权估计是多少？
\item
  为什么必须避免 \(p_{t,i}=0\)？
\item
  EXP3的遗憾比较对象是每轮最优臂还是事后最优的固定臂？
\end{enumerate}

\subsection{本章习题解答}\label{ux672cux7ae0ux4e60ux9898ux89e3ux7b54-4}

\begin{enumerate}
\def\labelenumi{\arabic{enumi}.}
\tightlist
\item
  所选臂的估计为 \(0.5/0.1=5\)。
\item
  一方面无法再观察该臂；另一方面重要性加权中的除法无定义。
\item
  事后最优的固定臂。
\end{enumerate}

\begin{center}\rule{0.5\linewidth}{0.5pt}\end{center}

\fi
\subsection{本章关键疑问与解答}\label{ux672cux7ae0ux5173ux952eux7591ux95ee-4}

\textbf{疑问一：对抗环境中的反馈可能剧烈变化，依据一次反馈更新权重还有意义吗？}

\textbf{解答：}单次反馈不是对某个臂永久好坏的判决，而是累计损失估计中的一次增量。EXP3 不假设某臂存在固定均值；它保证的是长期累计表现接近事后最优的固定臂。某次正反馈之后出现负反馈并不矛盾，后续更新会继续修正概率。

\textbf{疑问二：重要性加权既然无偏，为什么仍然需要控制最小抽样概率？}

\textbf{解答：}无偏只说明期望正确，不说明波动小。当 \(p_{t,i}\) 很小时，臂一旦被抽中，估计量会被除以很小的概率，从而产生巨大方差；若概率为零，甚至无法定义估计，也永远无法再获得该臂信息。EXP3 的探索混合和后文的 barrier、截断处理，本质上都在处理“无偏但不稳定”的问题。

\textbf{疑问三：为什么比较对象是事后最优固定臂，而不是每轮最优臂？}

\textbf{解答：}在完全对抗的序列中，每轮最优臂可能任意切换，bandit 反馈又只显示一个坐标；没有额外结构时，学习者无法追踪这样的动态基准。固定臂基准衡量的是：即使事后知道整段序列，始终选同一个臂最多能好多少。若要比较动态基准，必须加入切换次数、变差预算等额外限制。

\begin{center}\rule{0.5\linewidth}{0.5pt}\end{center}

\section{第六章
重尾分布与普通样本均值的失效}\label{ux7b2cux516dux7ae0-ux91cdux5c3eux5206ux5e03ux4e0eux666eux901aux6837ux672cux5747ux503cux7684ux5931ux6548}

\subsection{\texorpdfstring{6.1 有限 \(1+\epsilon\)
阶矩}{6.1 有限 1+\textbackslash epsilon 阶矩}}\label{ux6709ux9650-1epsilon-ux9636ux77e9}

\subsubsection{定义
6.1（重尾矩条件）}\label{ux5b9aux4e49-6.1ux91cdux5c3eux77e9ux6761ux4ef6}

设 \(\epsilon\in(0,1]\)。若

\[
\mathbb E|X|^{1+\epsilon}\le u,
\]

则称随机变量满足有限 \(1+\epsilon\) 阶原始矩条件。

\begin{itemize}
\tightlist
\item
  \(\epsilon=1\)：二阶矩有限，从而方差有限；
\item
  \(\epsilon=0.5\)：只要求 \(1.5\) 阶矩有限，方差可能无限；
\item
  \(\epsilon\) 越小，条件越宽松，允许尾部越重。
\end{itemize}

\subsection{6.2
为什么极端值会破坏普通均值？}\label{ux4e3aux4ec0ux4e48ux6781ux7aefux503cux4f1aux7834ux574fux666eux901aux5747ux503c}

考虑数据

\[
1,2,2,3,1000.
\]

普通均值为
\(201.6\)，远离大多数样本所在的区域。样本均值对每个观测赋予同样权重，因此单个极端值可以产生巨大影响。

在仅有 \(1+\epsilon\) 阶矩时，样本均值的偏差概率通常只能得到关于
\(1/\delta\)
的多项式依赖；UCB跨越多个时间点需要极小且可求和的失败概率，因此普通均值不再理想。

\subsection{6.3 截断均值}\label{ux622aux65adux5747ux503c}

\subsubsection{定义
6.2（截断均值）}\label{ux5b9aux4e49-6.2ux622aux65adux5747ux503c}

给定阈值 \(B_t>0\)，定义

\[
\widehat\mu_T
=\frac1n\sum_{t=1}^nX_t\mathbf1_{\{|X_t|\le B_t\}}.
\]

截断把原本可能无界的变量变成有界变量，使 Bernstein
型集中不等式可以使用。但截断会删除真实分布的一部分，因此会引入偏差。

\subsubsection{引理
6.1（截断偏差上界）}\label{ux5f15ux7406-6.1ux622aux65adux504fux5deeux4e0aux754c}

若 \(\mathbb E|X|^{1+\epsilon}\le u\)，则

\[
\mathbb E\left[|X|\mathbf1_{\{|X|>B\}}\right]
\le\frac{u}{B^\epsilon}.
\]

\subsubsection{证明}\label{ux8bc1ux660e-7}

在事件 \(|X|>B\) 上，\(|X|^\epsilon>B^\epsilon\)，因此

\[
|X|
=\frac{|X|^{1+\epsilon}}{|X|^\epsilon}
\le\frac{|X|^{1+\epsilon}}{B^\epsilon}.
\]

乘以示性函数并取期望：

\[
\mathbb E[|X|\mathbf1_{\{|X|>B\}}]
\le\frac{\mathbb E|X|^{1+\epsilon}}{B^\epsilon}
\le\frac{u}{B^\epsilon}.
\]

证毕。\(\square\)

\subsubsection{引理
6.2（截断后二阶矩上界）}\label{ux5f15ux7406-6.2ux622aux65adux540eux4e8cux9636ux77e9ux4e0aux754c}

\[
\mathbb E[X^2\mathbf1_{\{|X|\le B\}}]
\le uB^{1-\epsilon}.
\]

\subsubsection{证明}\label{ux8bc1ux660e-8}

在 \(|X|\le B\) 上，

\[
X^2=|X|^{1+\epsilon}|X|^{1-\epsilon}
\le |X|^{1+\epsilon}B^{1-\epsilon}.
\]

取期望并利用矩条件即得。\(\square\)

阈值选择体现偏差---波动权衡：

\begin{itemize}
\tightlist
\item
  \(B\) 太小：删除样本太多，截断偏差大；
\item
  \(B\) 太大：保留的极端值太多，随机波动大。
\end{itemize}

平衡两者后，可得到误差速度

\[
|\widehat\mu_T-\mu|
\lesssim
u^{1/(1+\epsilon)}
\left(\frac{\log(1/\delta)}n\right)^{\epsilon/(1+\epsilon)}.
\]

\subsection{6.4 Median-of-Means}\label{median-of-means}

\subsubsection{定义
6.3（Median-of-Means）}\label{ux5b9aux4e49-6.3median-of-means}

将 \(n\) 个样本分成 \(k\) 个互不相交的块，计算每个块的样本均值

\[
\widehat\mu_1,\ldots,\widehat\mu_k,
\]

再取这些块均值的中位数：

\[
\widehat\mu_{\mathrm{MoM}}
=\operatorname{median}(\widehat\mu_1,\ldots,\widehat\mu_k).
\]

极端值可能破坏少数块，但只要超过一半的块是好块，坏块就无法把中位数任意拖走。

\subsubsection{原始矩与中心矩}\label{ux539fux59cbux77e9ux4e0eux4e2dux5fc3ux77e9}

截断均值常依赖

\[
\mathbb E|X|^{1+\epsilon},
\]

而 Median-of-Means 可在中心矩

\[
\mathbb E|X-\mu|^{1+\epsilon}
\]

条件下工作。中心矩具有平移不变性：若 \(X'=X+a\)，则 \(\mu'=\mu+a\)，从而

\[
X'-\mu'=X-\mu.
\]

\iffalse

\begin{enumerate}
\def\labelenumi{\arabic{enumi}.}
\tightlist
\item
  若 \(u=4,\epsilon=1,B=2\)，计算截断偏差上界 \(u/B^\epsilon\)。
\item
  五个块均值为 \(2,2.1,1.9,500,2.2\)，求 Median-of-Means。
\item
  解释为什么对所有奖励同时加1000不会改变中心矩。
\end{enumerate}

\subsection{本章习题解答}\label{ux672cux7ae0ux4e60ux9898ux89e3ux7b54-5}

\begin{enumerate}
\def\labelenumi{\arabic{enumi}.}
\tightlist
\item
  上界为 \(4/2=2\)。
\item
  排序为 \(1.9,2,2.1,2.2,500\)，中位数为 \(2.1\)。
\item
  因为 \((X+1000)-(\mu+1000)=X-\mu\)。
\end{enumerate}

\begin{center}\rule{0.5\linewidth}{0.5pt}\end{center}

\fi
\subsection{本章关键疑问与解答}\label{ux672cux7ae0ux5173ux952eux7591ux95ee-5}

\textbf{疑问一：普通样本均值无偏，为什么在重尾环境中仍可能不是合适的估计器？}

\textbf{解答：}无偏只描述平均意义，不保证高概率误差小。重尾分布会以小概率产生极端样本，一个极端值就可能长时间主导样本均值；仅有限低阶矩时，经典指数集中界不再成立。bandit 算法需要的是可转化为置信区间的高概率控制，而不只是期望正确。

\textbf{疑问二：截断为什么会同时带来好处和偏差？阈值应怎样理解？}

\textbf{解答：}截断限制单个样本的影响，使处理后的随机变量能够使用更强的集中工具；但被删除或压缩的尾部带来系统偏差。阈值小则波动容易控制但偏差大，阈值大则偏差小但极端值影响增强。阈值随样本量增长，是为了让早期稳健性和长期一致性同时成立。

\textbf{疑问三：Median-of-Means 为什么能抵抗少数极端值，它与普通中位数有何不同？}

\textbf{解答：}它先分块取均值，让每个块利用平均化降低普通波动，再对块均值取中位数，使少数被极端值污染的坏块无法支配结果。它估计的是总体均值，不是总体中位数；稳健性来自“多数块是好块”，而不是把原始观测直接排序。

\begin{center}\rule{0.5\linewidth}{0.5pt}\end{center}

\section{第七章 Robust UCB
与重尾遗憾界}\label{ux7b2cux4e03ux7ae0-robust-ucb-ux4e0eux91cdux5c3eux9057ux61beux754c}

\subsection{7.1
稳健均值估计器接口}\label{ux7a33ux5065ux5747ux503cux4f30ux8ba1ux5668ux63a5ux53e3}

假设存在估计量 \(\widehat\mu(n,\delta)\)，使得以至少 \(1-\delta\)
的概率，

\[
|\widehat\mu-\mu|
\le
v^{1/(1+\epsilon)}
\left(\frac{c\log(1/\delta)}n\right)^{\epsilon/(1+\epsilon)}.
\]

这一条件可以看作算法接口：只要一个均值估计器满足它，就可以被装入 Robust
UCB。截断均值、Median-of-Means 和 Catoni
估计器分别在不同矩条件下提供这种保证。

\subsection{7.2 Robust UCB 指标}\label{robust-ucb-ux6307ux6807}

若臂 \(i\) 已被选择 \(s\) 次，定义

\[
B_{i,s,t}
=\widehat\mu_{i,s,t}
+v^{1/(1+\epsilon)}
\left(\frac{c\log(t^2)}s\right)^{\epsilon/(1+\epsilon)}.
\]

第 \(t\) 轮选择指标最大的臂。这里相当于取
\(\delta=t^{-2}\)，使各轮失败概率可以求和控制。

当 \(\epsilon=1\) 时，置信半径按 \(s^{-1/2}\) 收缩；当 \(\epsilon=0.5\)
时，它只按 \(s^{-1/3}\) 收缩。因此尾部越重，均值越难估计。

\subsection{7.3 主定理}\label{ux4e3bux5b9aux7406}

\subsubsection{定理 7.1（Robust UCB
遗憾界，忽略常数）}\label{ux5b9aux7406-7.1robust-ucb-ux9057ux61beux754cux5ffdux7565ux5e38ux6570}

若各臂满足有限 \(1+\epsilon\)
阶矩条件，且所用稳健估计器满足上述置信界，则

\[
R_T
=O\left(
\sum_{i:\Delta_i>0}
\frac{v^{1/\epsilon}\log T}{\Delta_i^{1/\epsilon}}
\right).
\]

\subsubsection{证明骨架}\label{ux8bc1ux660eux9aa8ux67b6}

对次优臂 \(i\)，如果置信区间均正确，那么只要

\[
2v^{1/(1+\epsilon)}
\left(\frac{c\log T}{N_i}\right)^{\epsilon/(1+\epsilon)}
\le\Delta_i,
\]

该臂就不会继续优先于最优臂。解出 \(N_i\)：

\[
N_i
=O\left(
\frac{v^{1/\epsilon}\log T}
{\Delta_i^{(1+\epsilon)/\epsilon}}
\right).
\]

置信失败事件的概率在时间上可求和，只贡献常数阶的额外选择次数。因此

\[
\mathbb E[T_i(T)]
=O\left(
\frac{v^{1/\epsilon}\log T}
{\Delta_i^{(1+\epsilon)/\epsilon}}
\right).
\]

最后乘以每次选择臂 \(i\) 的损失 \(\Delta_i\)：

\[
\begin{aligned}
R_{T,i}
&=\Delta_i\mathbb E[T_i(T)]\\
&=O\left(
\frac{v^{1/\epsilon}\log T}
{\Delta_i^{1/\epsilon}}
\right).
\end{aligned}
\]

对所有次优臂求和即得结论。\(\square\)

\subsection{7.4
如何解释定理？}\label{ux5982ux4f55ux89e3ux91caux5b9aux7406}

当 \(\epsilon=1\) 时，

\[
R_T=O\left(\sum_i\frac{v\log T}{\Delta_i}\right),
\]

与 sub-Gaussian UCB
的阶数相同。因此有限方差已经足以恢复经典遗憾阶数，前提是不用普通均值，而使用稳健估计器。

当 \(\epsilon<1\) 时，对小 gap 的依赖恶化。例如 \(\epsilon=0.5\) 时，

\[
R_T=O\left(\sum_i\frac{v^2\log T}{\Delta_i^2}\right).
\]

这说明重尾使相近臂更难区分。相应下界表明 \(1/\Delta_i^{1/\epsilon}\)
的依赖在一般情形下无法消除。

\iffalse

\begin{enumerate}
\def\labelenumi{\arabic{enumi}.}
\tightlist
\item
  当 \(\epsilon=0.5\) 时，置信半径关于样本数按什么速度收缩？
\item
  当 \(\Delta=0.1\) 时，忽略其他项，比较 \(\epsilon=1\) 和
  \(\epsilon=0.5\) 时的 gap 因子。
\item
  解释``置信区间都正确且次优臂已经采样充分时，算法不会继续大量选择它''。
\end{enumerate}

\subsection{本章习题解答}\label{ux672cux7ae0ux4e60ux9898ux89e3ux7b54-6}

\begin{enumerate}
\def\labelenumi{\arabic{enumi}.}
\tightlist
\item
  \(\epsilon/(1+\epsilon)=1/3\)，故按 \(n^{-1/3}\) 收缩。
\item
  \(\epsilon=1\) 时为 \(1/0.1=10\)；\(\epsilon=0.5\) 时为
  \(1/0.1^2=100\)。后者更难。
\item
  此时次优臂的UCB不超过 \(\mu_i+\Delta_i=\mu^*\)，而最优臂UCB至少为
  \(\mu^*\)，故次优臂不再严格占优。
\end{enumerate}

\begin{center}\rule{0.5\linewidth}{0.5pt}\end{center}

\fi
\subsection{本章关键疑问与解答}\label{ux672cux7ae0ux5173ux952eux7591ux95ee-6}

\textbf{疑问一：Robust UCB 与经典 UCB 的证明骨架为什么几乎相同，结果却对 gap 有更差的依赖？}

\textbf{解答：}两者都依靠“置信区间正确且半径小于 gap 的一半后，次优臂不再占优”。差别在置信半径的收缩速度：重尾条件下半径约按 \(n^{-\epsilon/(1+\epsilon)}\) 收缩，比 \(n^{-1/2}\) 慢。为把半径压到 \(\Delta_i\) 以下需要更多样本，因此遗憾中的 gap 幂次恶化。

\textbf{疑问二：为什么把稳健均值估计器写成一个“接口”是有价值的？}

\textbf{解答：}UCB 的决策逻辑只需要估计值及其高概率误差界，并不依赖估计器内部究竟使用截断、Median-of-Means 还是 Catoni 方法。先证明抽象接口下的遗憾结论，再代入不同估计器，可以把概率估计问题与序贯决策问题分离，也能清楚看出每种矩条件影响的是哪一步。

\textbf{疑问三：有限方差为何可能恢复经典的 \(\log T/\Delta_i\) 遗憾阶数，却仍不能直接使用普通均值的次高斯置信区间？}

\textbf{解答：}有限方差允许构造具有次高斯型高概率误差的稳健估计器，但普通样本均值本身在所有有限方差分布上并没有同样的指数尾保证。恢复经典阶数靠的是更合适的估计方法，而不是把“有限方差”误认为“次高斯”。

\begin{center}\rule{0.5\linewidth}{0.5pt}\end{center}

\section{第八章 从经典重尾老虎机到
uniINF}\label{ux7b2cux516bux7ae0-ux4eceux7ecfux5178ux91cdux5c3eux8001ux864eux673aux5230-uniinf}

\subsection{8.1
经典方法的局限}\label{ux7ecfux5178ux65b9ux6cd5ux7684ux5c40ux9650}

经典 Robust UCB
主要面向随机平稳环境，并且常要求预先知道重尾阶数和矩上界。现实中，学习者可能不知道：

\begin{enumerate}
\def\labelenumi{\arabic{enumi}.}
\tightlist
\item
  环境究竟是随机平稳还是对抗变化的；
\item
  重尾阶数 \(\alpha\) 或 \(1+\epsilon\)；
\item
  尾部尺度 \(\sigma\) 或矩上界 \(u\)。
\end{enumerate}

\subsection{8.2 Best-of-Both-Worlds}\label{best-of-both-worlds}

一个 Best-of-Both-Worlds
算法使用同一套更新规则，在不知道环境类型的情况下：

\begin{itemize}
\tightlist
\item
  在随机环境中取得较小的 gap-dependent 遗憾；
\item
  在对抗环境中取得接近最优的 worst-case 遗憾。
\end{itemize}

\subsection{8.3 Parameter-Free}\label{parameter-free}

这里的 parameter-free
不是``算法没有任何设置''，而是不需要预先知道未知的重尾参数。算法需要根据实际观测自动调节学习率与截断尺度。

\subsection{8.4
为什么重尾和重要性加权结合后更困难？}\label{ux4e3aux4ec0ux4e48ux91cdux5c3eux548cux91cdux8981ux6027ux52a0ux6743ux7ed3ux5408ux540eux66f4ux56f0ux96be}

对抗型老虎机使用

\[
\widehat\ell_{t,i}
=\frac{\ell_{t,i}\mathbf1_{\{A_t=i\}}}{p_{t,i}}.
\]

当损失 \(\ell_{t,i}\) 本身可能是重尾极端值，同时选择概率 \(p_{t,i}\)
很小时，商

\[
\frac{\ell_{t,i}}{p_{t,i}}
\]

可能异常巨大。因此需要自适应
clipping/skipping、概率正则化和学习率调度共同控制估计量。


\clearpage
\subsection{8.6 阅读定位、原文与符号约定}

本节起为补充的 uniINF 研读内容。依据：Yu Chen、Jiatai Huang、Yan Dai、Longbo Huang，\emph{uniINF: Best-of-Both-Worlds Algorithm for Parameter-Free Heavy-Tailed MABs}，ICLR 2025；使用所提供的 arXiv:2410.03284v2（2025-02-25，共 37 页）版本。原文链接：\url{https://arxiv.org/abs/2410.03284v2}。

\textbf{截至目前的学习状态。}
\begin{itemize}
\item \textbf{能够独立说明并完成基础推导：}重尾矩条件、损失版伪遗憾、重要性加权的条件无偏性、FTRL 的基本定义、负熵产生指数权重、log-barrier 的梯度与 KKT 方程、尾部一阶矩和截断二阶矩界，以及随机环境中的 self-bounding 基础关系。
\item \textbf{已经理解算法结构，但仍需通过复述和手算巩固：}uniINF 中 FTRL、skipping、clipping 和自适应尺度的分工，parameter-free 与 Best-of-Both-Worlds 的含义，以及主定理对 $K,T,\alpha,\sigma,\Delta_{\min}$ 的依赖。
\item \textbf{仍在研读：}概率乘法稳定性的完整证明、精细 Bregman 散度界、$\Psi$-shifting 全部细节、skipping 误差的停止时刻分析，以及附录中的完整主定理证明。
\end{itemize}

\textbf{阅读顺序。}先读原文第 4 页 Section 3，再读第 5 页 Algorithm 1 和第 5--6 页 Section 4；之后看第 7 页 Theorem 3、第 7--11 页证明框架。附录 C--F 留作第二轮精读。

\begin{itemize}
\item 原文用 $i_t$ 表示所选臂，本笔记统一改用 $A_t$；原文概率向量 $x_t$ 对应前文的 $p_t$，下文沿用 $x_t$ 方便对照。
\item 改用\textbf{损失最小化}：$\ell_{t,i}$ 越小越好。随机环境中 $\mu_i=\mathbb E[\ell_{t,i}]$，$\Delta_i=\mu_i-\mu_{i^*}$；注意与奖励最大化时 gap 的方向相反。
\item $\mathcal F_{t-1}$ 表示第 $t$ 轮行动前的历史信息。$x_t,S_t,C_{t,i}$ 是由历史决定的随机量，但给定 $\mathcal F_{t-1}$ 后可视为已知。
\item $A_t,\ell_{t,i},N_i(T)$ 是随机变量。固定环境下的 $\mu_i,\Delta_i$ 及下述已取期望的 $R_T$ 是确定量。实际累计损失差不等于其期望。
\end{itemize}

\subsection{8.7 设定：}

\subsubsection{定义 8.1（重尾损失模型与伪遗憾）}
每轮先根据历史选取概率向量 $x_t\in\Delta_K$，再抽样 $A_t\sim x_t$，只观察 $\ell_{t,A_t}$。其中
\[
\Delta_K=\{x\in\mathbb R^K:x_i\ge0,\ \sum_i x_i=1\},\qquad
R_T=\max_{i\in[K]}\mathbb E\!\left[\sum_{t=1}^T(\ell_{t,A_t}-\ell_{t,i})\right].
\]
这里的比较对象是一个\textbf{固定臂}，不是每轮事后选出的最好臂；$\max_i\mathbb E[\cdot]$ 也不能随意换成 $\mathbb E[\max_i\cdot]$。

原文假定存在未知的 $\alpha\in(1,2]$ 和 $\sigma\ge0$，使所有轮次与臂均满足
\[
\mathbb E|\ell_{t,i}|^\alpha\le\sigma^\alpha.
\]
这与前文 $1+\epsilon$ 阶矩的记号对应为 $\alpha=1+\epsilon$。$\alpha=2$ 给出有限二阶原始矩，不代表次高斯；$1<\alpha<2$ 则\emph{允许}方差不存在，但不强制方差无穷。

随机环境的分布不随时间变化；原文对抗环境允许分布随 $t,i$ 变化，但不依赖学习者之前的行动，即 oblivious adversary。不能把论文结论直接推广到看见当前行动后才定损失的对手。

\subsubsection{假设 8.1（最优臂的截断非负性，原文 Assumption 1）}
存在最优臂 $i^*$，对每个 $t$ 和所有 $M\ge0$，有
\[
\mathbb E[\ell_{t,i^*}\mathbf1_{\{|\ell_{t,i^*}|>M\}}]\ge0.
\]
它限制的是最优臂尾部的\textbf{带符号期望}，不是要求每次损失非负，也不是要求每个臂都非负。例如，非负随机变量一定满足；可积且关于零对称的随机变量也满足，因为两侧尾部贡献抵消。仅知道 $\mathbb EX\ge0$ 不够，见习题。

\textbf{随机情形的额外假设（原文 Assumption 2）。}最优臂唯一，即
\[
\Delta_{\min}=\min_{i\ne i^*}\Delta_i>0.
\]
因此，“Parameter-Free”指算法不需要输入 $\alpha,\sigma$，不是“对任意分布都有同样保证”，也不是无需 $K,T$、初始尺度等任何设置。

\subsection{8.8 从 EXP3 到 FTRL：先理解优化问题}

\subsubsection{定义 8.2（FTRL）}
令 $L_{t-1}=\sum_{s<t}\widetilde\ell_s$ 为累计估计损失。跟随正则化领导者选择
\[
x_t=\arg\min_{x\in\Delta_K}\{\langle L_{t-1},x\rangle+\Psi_t(x)\}.
\]
第一项追求过去损失小，第二项调节概率分配，避免只因短期估计就把概率压到边界。FTRL 是\textbf{框架}；换正则项、估计量或尺度调度，得到的算法及保证都会不同。UCB 依据置信上界选臂，并不是此处的 FTRL 更新。

\subsubsection{命题 8.1（负熵产生指数权重）}
固定 $\eta>0$，若 $\Psi(x)=\eta^{-1}\sum_i x_i\log x_i$，则
\[
x_i=\frac{\exp(-\eta L_i)}{\sum_j\exp(-\eta L_j)}.
\]
\textbf{证明。}对约束 $\sum_i x_i=1$ 引入乘子 $\lambda$，一阶条件为
\[
L_i+\eta^{-1}(1+\log x_i)+\lambda=0.
\]
于是 $x_i=\exp(-1-\eta\lambda)\exp(-\eta L_i)$，由归一化得到结论。目标严格凸，正的驻点即唯一最优解。证毕。

这解释了 EXP3 指数权重更新与 FTRL 的联系；具体 EXP3 版本可能另有显式探索混合和不同学习率，不能把所有版本不加区分地写成同一个公式。

\subsubsection{定义 8.3（uniINF 的 log-barrier）}
uniINF 取
\[
\Psi_t(x)=-S_t\sum_{i=1}^K\log x_i,\qquad S_1=4.
\]
定义域是单纯形内部 $\Delta_K^\circ=\{x_i>0,\sum_i x_i=1\}$；边界上视为 $+\infty$。其梯度与 Hessian 为
\[
(\nabla\Psi_t(x))_i=-S_t/x_i,\qquad
\nabla^2\Psi_t(x)=S_t\operatorname{diag}(x_i^{-2})\succ0.
\]
当某个 $x_i\downarrow0$ 时惩罚趋于无穷，因此每个臂概率保持正值；这\textbf{不等于}存在一个与时间无关的统一正下界。

\subsubsection{命题 8.2（KKT 方程化为一维求根）}
固定有限向量 $L$ 与 $S>0$，最优解满足
\[
x_i=\frac{S}{L_i+\lambda},\qquad
\sum_{i=1}^K\frac{S}{L_i+\lambda}=1,\qquad \lambda>-\min_i L_i.
\]
\textbf{证明。}目标严格凸且在边界发散，故存在唯一内部最优解。一阶条件 $L_i-S/x_i+\lambda=0$ 给出上述形式。设 $g(\lambda)=\sum_iS/(L_i+\lambda)$，则
\[
g'(\lambda)=-\sum_i\frac{S}{(L_i+\lambda)^2}<0.
\]
在左端 $g\to+\infty$，在 $\lambda\to+\infty$ 时 $g\to0$。由连续性和严格单调性，$g(\lambda)=1$ 恰有一根。证毕。

\textbf{解释。}原文附录写作 $x_{t,i}=S_t/(L_{t-1,i}-Z_t)$，这里只是 $\lambda=-Z_t$。若 $L_i<L_j$，则 $x_i>x_j$；但概率不是指数形式。$S_t$ 是正则化尺度，通常把 $1/S_t$ 看作有效学习率：$S_t$ 增大不表示更新更激进。

\subsection{8.9 uniINF 在做什么？}

下面按原文 Algorithm 1（第 5 页，公式 (2)--(4)）重写，假定 $K\ge2,T\ge2$。单臂问题本身无须探索。$T$ 是算法中的给定轮数。

\begin{enumerate}
\item 用 $L_{t-1}$ 和 $S_t$ 解上一节 FTRL 问题，得到 $x_t$。
\item 抽样 $A_t\sim x_t$，仅观察 $\ell_{t,A_t}$。
\item 对所选臂计算阈值 $C_{t,A_t}=S_t/[4(1-x_{t,A_t})]$。理论上对每个臂定义
\[
C_{t,i}=\frac{S_t}{4(1-x_{t,i})},\quad
\ell^{\rm skip}_{t,i}=\ell_{t,i}\mathbf1_{\{|\ell_{t,i}|<C_{t,i}\}},\quad
\ell^{\rm clip}_{t,i}=\max\{-C_{t,i},\min(\ell_{t,i},C_{t,i})\}.
\]
\item 更新累计估计损失：
\[
\widetilde\ell_{t,i}=\frac{\ell^{\rm skip}_{t,i}}{x_{t,i}}\mathbf1_{\{A_t=i\}},\qquad
L_t=L_{t-1}+\widetilde\ell_t.
\]
\item 用 clipped 损失更新尺度：
\[
\boxed{S_{t+1}^2=S_t^2+
\frac{(\ell^{\rm clip}_{t,A_t})^2(1-x_{t,A_t})^2}{K\log T}.}
\]
\end{enumerate}

\textbf{不能混淆的两条路径。}Skipped 损失进入累计估计 $L_t$，从而影响后续 FTRL；clipped 损失进入 $S_{t+1}$，从而改变正则化尺度和后续阈值。理论中定义了整条损失向量，不表示算法能观测所有臂的反馈。

\textbf{Skipping 不会抹去真实损失。}即使本轮 $\widetilde\ell_t=0$，实际付出的仍是 $\ell_{t,A_t}$，并非零。遗憾分析必须为这部分差异付出代价。

\subsubsection{例 8.1（从零状态完整手算一轮）}
取 $K=2,T=100,S_1=4,L_0=(0,0)$。对称性给出 $x_1=(1/2,1/2)$，故 $C_{1,1}=C_{1,2}=2$。假设抽到臂 1，观察损失 $10$，则
\[
\ell^{\rm skip}_{1,1}=0,\quad \ell^{\rm clip}_{1,1}=2,\quad
\widetilde\ell_1=(0,0),\quad L_1=(0,0),
\]
\[
S_2=\sqrt{16+\frac{1}{2\log100}}\approx4.01355.
\]
下一轮仍因累计损失对称而均匀抽样，但阈值变为 $S_2/2\approx2.00677$。本轮真实损失是 $10$；尺度仍然有所增长。如果误用 skipped 损失更新尺度，则 $S_2=4$，反复遇到大损失时可能一直停在原阈值。

\subsubsection{命题 8.3（发生 skipping 时，尺度仍按比例增长）}
如果 $|\ell_{t,A_t}|\ge C_{t,A_t}$，则
\[
S_{t+1}=S_t\sqrt{1+\frac{1}{16K\log T}}.
\]
\textbf{证明。}这时 $(\ell^{\rm clip}_{t,A_t})^2=C_{t,A_t}^2$，代入阈值可得
$C_{t,A_t}^2(1-x_{t,A_t})^2=S_t^2/16$，再代入尺度更新即可。证毕。

\subsection{8.10 偏差从哪里来，有限矩如何控制它？}

\subsubsection{命题 8.4（重要性加权的条件无偏性）}
在原文环境设定下，给定 $\mathcal F_{t-1}$ 和当轮损失向量 $\ell_t$，抽样概率为 $x_t$，因此
\[
\mathbb E[\widetilde\ell_{t,i}\mid\mathcal F_{t-1},\ell_t]
=\frac{\ell^{\rm skip}_{t,i}}{x_{t,i}}\mathbb P(A_t=i\mid\mathcal F_{t-1},\ell_t)
=\ell^{\rm skip}_{t,i}.
\]
再取条件期望，得到
$\mathbb E[\widetilde\ell_{t,i}\mid\mathcal F_{t-1}]=\mathbb E[\ell^{\rm skip}_{t,i}\mid\mathcal F_{t-1}]$。证毕。

它对\textbf{skipped 损失}无偏，通常不对原损失无偏。差值为被跳过的尾部期望。

\subsubsection{命题 8.5（尾部一阶矩与截断二阶矩）}
若 $\mathbb E|X|^\alpha\le\sigma^\alpha$、$1<\alpha\le2$、$C>0$，则
\[
\mathbb E[|X|\mathbf1_{\{|X|\ge C\}}]\le\sigma^\alpha C^{1-\alpha},\qquad
\mathbb E[\operatorname{Clip}(X,C)^2]\le\sigma^\alpha C^{2-\alpha}.
\]
\textbf{证明。}在 $|X|\ge C$ 上，$|X|\le |X|^\alpha C^{1-\alpha}$，取期望得第一式。第二式用逐点不等式
\[
\min(|X|,C)^2\le |X|^\alpha C^{2-\alpha}.
\]
若 $|X|\le C$，由 $|X|^{2-\alpha}\le C^{2-\alpha}$ 得到；若 $|X|>C$，由 $C^\alpha\le|X|^\alpha$ 得到。取期望即证。证毕。

第一式控制 skipping 的误差；第二式在原损失二阶矩可能无穷时，仍控制处理后估计的二阶项。阈值越大，尾部误差界越小，但二阶项的界通常越大。这是阈值设计要平衡的两种代价。

\subsection{8.11 Bregman 散度：概率变化的分析工具}

\subsubsection{定义 8.4 与命题 8.6（非负性和具体形式）}
对可微凸函数 $\Psi$，定义
\[
D_\Psi(u,v)=\Psi(u)-\Psi(v)-\langle\nabla\Psi(v),u-v\rangle.
\]
凸函数的一阶支撑不等式直接给出 $D_\Psi(u,v)\ge0$。它一般不对称，也未必满足三角不等式，不能当成通常的距离。

对 $\Psi(x)=-S\sum_i\log x_i$，逐项代入可得
\[
D_\Psi(u,v)=S\sum_i\left[-\log\frac{u_i}{v_i}+\frac{u_i}{v_i}-1\right].
\]
单项非负也可由 $\log q\le q-1$（$q>0$）直接验证。

原文分析中的 $z_t$ 是加入当轮估计后、\textbf{仍保持尺度 $S_t$} 的最优解：
\[
z_t=\arg\min_z\{\langle L_t,z\rangle+\Psi_t(z)\}.
\]
真正的 $x_{t+1}$ 使用 $S_{t+1}$，所以一般 $z_t\ne x_{t+1}$。令 $\mathrm{DIV}_t=D_{\Psi_t}(x_t,z_t)$。

\subsubsection{论文技术引用}
原文 Lemmas 4--5 建立概率的乘法稳定性，并给出如下关键量级：
\[
\frac{x_{t,i}}2\le z_{t,i}\le2x_{t,i},\qquad
\mathrm{DIV}_t\lesssim S_t^{-1}(\ell^{\rm skip}_{t,A_t})^2(1-x_{t,A_t})^2.
\]

\textbf{一个应自行核查的小步骤。}由 $x_i/2\le z_i\le2x_i$，只能推出 $w_i=x_i/z_i-1\in[-1/2,1]$：
\[
w-\log(1+w)\le w^2,\qquad w\ge-1/2.
\]
证明：设 $h(w)=w^2-w+\log(1+w)$，则
$h'(w)=w(1+2w)/(1+w)$。在 $[-1/2,0]$ 上 $h'\le0$，在 $[0,\infty)$ 上 $h'\ge0$，故 $h(w)\ge h(0)=0$。

\subsection{8.12 为什么随机环境能得到更小的遗憾？}

\subsubsection{命题 8.7（self-bounding 的基础关系）}
在随机环境下，
\[
R_T=\sum_{t=1}^T\sum_{i\ne i^*}\Delta_i\mathbb E[x_{t,i}]
\ge\Delta_{\min}\sum_{t=1}^T\mathbb E[1-x_{t,i^*}].
\]
\textbf{证明。}给定历史，当轮平均损失差为 $\sum_i x_{t,i}(\mu_i-\mu_{i^*})$。取期望并求和得等式；再用 $\Delta_i\ge\Delta_{\min}$ 及 $\sum_{i\ne i^*}x_{t,i}=1-x_{t,i^*}$ 得不等式。证毕。

\subsubsection{从二阶项到最优臂的剩余概率}
把命题 8.5 与阈值代入技术引理，可得某常数 $c_0>0$ 下
\begin{align*}
\mathbb E[\mathrm{DIV}_t\mid\mathcal F_{t-1}]
&\le c_0S_t^{-1}\sum_i x_{t,i}\sigma^\alpha C_{t,i}^{2-\alpha}(1-x_{t,i})^2\\
&=c_0 4^{\alpha-2}\sigma^\alpha S_t^{1-\alpha}\sum_i x_{t,i}(1-x_{t,i})^\alpha\\
&\le 2c_0\sigma^\alpha S_t^{1-\alpha}(1-x_{t,i^*}).
\end{align*}
最后一步因为非最优臂的和不超过 $1-x_{t,i^*}$；最优臂项也不超过 $1-x_{t,i^*}$。于是技术引理中的 $(1-x)^2$ 最终连接到 self-bounding 关系，而非仅仅表示“多探索”。

\subsubsection{停止阈值与遗憾吸收（证明思路）}
将上式常数合并为 $c$。希望当 $S_t\ge M$ 时有 $c\sigma^\alpha S_t^{1-\alpha}\le\Delta_{\min}/4$，只需
\[
M\ge\left(\frac{4c\sigma^\alpha}{\Delta_{\min}}\right)^{1/(\alpha-1)}.
\]
由于 $1-\alpha<0$，$S_t$ 越大，这个上界越小。令
\[
\tau=\min\bigl(\{t\in\{1,\ldots,T\}:S_t\ge M\}\cup\{T+1\}\bigr).
\]
达到阈值与否由已有信息确定，因此能把求和按 $t<\tau$ 和 $t\ge\tau$ 划分；这不是随意套用“可选停止定理”。前段利用累计尺度界控制，后段利用 self-bounding 吸收进 $R_T/4$。

\textbf{重要：}$M$ 依赖未知参数，但只用于\emph{证明}，不是算法输入。分析者知道如何选择 $M$，并不意味着学习者需要知道它。

\subsection{8.13 遗憾分解与主结果应该怎样读}

\subsubsection{三项分解：先做代数，再使用假设}
令 $y=e_{i^*}$。log-barrier 在边界 $y$ 上为无穷，分析时需要内部比较点。为避免短轮数的边界问题，本节取
$\bar y=(1-1/T)y+(1/T)\mathbf1/K$（$T>1$）。原文采用另一种 $1/T$ 级别的平滑，思想相同。恒等式为
\begin{align*}
\langle x_t-y,\ell_t\rangle
={}&\langle\bar y-y,\ell_t^{\rm skip}\rangle
+\langle x_t-\bar y,\ell_t^{\rm skip}\rangle\\
&+\langle x_t-y,\ell_t-\ell_t^{\rm skip}\rangle.
\end{align*}
三项分别是比较点平滑误差、FTRL 主遗憾、skipping 误差。由 $\mathbb E|\ell_{t,i}|\le\sigma$，本节平滑点的累计误差绝对值不超过 $2\sigma$；原文平滑方式得到 $O(\sigma K)$ 级别项。

\textbf{截断非负性在哪里使用？}令 $d_{t,i}=\ell_{t,i}-\ell^{\rm skip}_{t,i}$。给定历史，最优臂项为
\[
(x_{t,i^*}-1)\mathbb E[d_{t,i^*}\mid\mathcal F_{t-1}]\le0,
\]
因为第一个因子非正，第二个由假设非负。因此 skipping 误差可以只上界为非最优臂尾部误差的和。这是\emph{条件期望}中的结论，不是每次实现值都满足。原文假设写 $>M$、算法使用 $\ge C$，对可积变量可由 $M\uparrow C$ 与支配收敛衔接。

FTRL 主遗憾进一步分成 $\mathrm{DIV}_t$ 与尺度变化引起的 $\mathrm{SHIFT}_t$。原文分别控制这两项和非最优臂 skipping 项；随机情形中三项可各贡献至多 $R_T/4$ 加一个可控余项。合起来若
\[
R_T\le B+\frac34R_T,
\]
移项便得到 $R_T\le4B$。真正困难在于逐项证明余项界，而不在最后移项。

\subsubsection{原文 Theorem 3：主导依赖关系}
在原文设定和 Assumption 1 下，对抗情形的主导量级为
\[
R_T=\widetilde O\bigl(\sigma K^{1-1/\alpha}T^{1/\alpha}\bigr).
\]
随机情形再加入 Assumption 2，其定理展示的主导关系为
\[
R_T=O\!\left(K\left(\frac{\sigma^\alpha}{\Delta_{\min}}\right)^{1/(\alpha-1)}
\log T\cdot\log\frac{\sigma^\alpha}{\Delta_{\min}}\right).
\]
\textbf{使用范围。}以上是原文的渐近量级表述，不作为对任意尺度都直接代入的非负数值上界。尤其当 $\sigma^\alpha/\Delta_{\min}\le1$ 时，不能将最后一个裸对数按负数使用；分析需保留初始尺度、常数项和适当正部对数。固定 $\alpha$ 时隐藏的常数也可能随 $\alpha\downarrow1$ 恶化。本笔记不声称修订了原文的全参数有限样本定理。

\begin{itemize}
\item $\alpha=2$：对抗主导项是 $\widetilde O(\sigma\sqrt{KT})$；随机项对 gap 的主要依赖是 $\sigma^2/\Delta_{\min}$。
\item $\alpha=3/2$：对抗主导项为 $\widetilde O(\sigma K^{1/3}T^{2/3})$；随机项的 gap 幂次变为 $\Delta_{\min}^{-2}$。
\item 固定参数时随机项对 $T$ 为对数量级；这来自同一个算法的环境适应性，而不是先知道环境类型再切换 UCB/EXP3。
\item uniINF 的随机界使用 $K$ 与 $\Delta_{\min}$，不能直接写成经典下界中的逐臂求和
\[
\sum_{i\ne i^*}(\sigma^\alpha/\Delta_i)^{1/(\alpha-1)}\log T.
\]
当各 gap 差异很大时，两者可能相差较多；原文 Table 1 与结论也讨论了这一局限。
\end{itemize}

\iffalse
\begin{enumerate}
\item 设 $X$ 以概率 $0.9$ 取 $1$、以概率 $0.1$ 取 $-5$。求 $\mathbb EX$，并判断它是否满足截断非负性。
\item 对 $L=(0,0),S=4,K=2$，求 log-barrier FTRL 的 $\lambda$ 与 $x$。若改为 $L_1<L_2$，哪个臂概率更大？
\item 取 $S_t=4,x_t=(0.8,0.2),K=2,T=100$。抽到臂 2 并观察损失 $10$，求阈值、skip、clip、估计向量和 $S_{t+1}$。若损失为 $-10$ 呢？
\item 在题 3 的状态下，若臂 2 的当轮损失恒为 $1$，证明其重要性加权估计的条件期望为 $1$。
\item 推导 $\Psi(x)=-S\sum_i\log x_i$ 的 Hessian，并解释“每个概率为正”是否意味着“每个概率至少为 $0.01$”。
\item 若 $\alpha=3/2$，将 $\sigma^\alpha M^{1-\alpha}\le\Delta_{\min}/4$ 解成对 $M$ 的要求；说明这个 $M$ 是否破坏 parameter-free。
\item 若 $R_T\le B+3R_T/4$，求 $R_T$ 的上界。说明仅写出这个式子是否等于证明主定理。
\item 不看公式，用四句话解释 uniINF 一轮的两条损失处理路径，并指出真实损失是否会因为 skipping 而消失。
\end{enumerate}

\subsection{8.15 习题解答与下一步精读}
\begin{enumerate}
\item $\mathbb EX=0.4>0$；取 $M=2$，有 $\mathbb E[X\mathbf1_{\{|X|>2\}}]=-0.5<0$。因此不满足。均值非负不足以控制每个截断阈值下的尾部。
\item $8/\lambda=1$，所以 $\lambda=8$、$x=(1/2,1/2)$。若 $L_1<L_2$，同一 $\lambda$ 下 $S/(L_1+\lambda)>S/(L_2+\lambda)$。
\item $C_{t,2}=4/(4\times0.8)=1.25$；skip 为 $0$，clip 为 $1.25$，估计向量为 $(0,0)$。尺度满足
\[
S_{t+1}=\sqrt{16+\frac{1.25^2\times0.8^2}{2\log100}}
=\sqrt{16+\frac1{2\log100}}\approx4.01355.
\]
若损失是 $-10$，clip 为 $-1.25$、skip 仍为零、尺度相同；真实损失的符号则不同。$x_t$ 是给定中间状态，不声称它来自 $L_0=0$。
\item 因 $1<1.25$，未被 skip。以 $0.2$ 概率得到估计 $1/0.2=5$，以 $0.8$ 概率得到 $0$，故均值 $0.2\times5=1$。未选臂的零是估计值，不是观察到其真实损失为零。
\item Hessian 为 $S\operatorname{diag}(x_i^{-2})$；正定支持严格凸性。每个概率为正不提供题目所述统一下界。
\item $\sigma^{3/2}/\sqrt M\le\Delta_{\min}/4$，所以 $M\ge16\sigma^3/\Delta_{\min}^2$。它只在证明中划分轮次，不传入算法，不破坏 parameter-free。
\item $R_T\le4B$。还必须证明三类误差都能被相应余项加 $R_T/4$ 控制，且说明假设的使用位置。
\item 先用累计估计损失和 log-barrier 求概率，再抽臂并观察损失；skip 后的重要性加权值进入累计损失；clip 后的平方值更新尺度；被 skip 的真实损失仍存在，因此需要分析 skipping 误差。
\end{enumerate}

\fi
\subsection{8.14 关键疑问与解答}

\textbf{疑问一：uniINF 为什么既要 skipping，又要 clipping，使用一种处理方式不够吗？}

\textbf{解答：}两者服务于不同的分析路径。Skipping 把超过阈值的损失从重要性加权估计中移除，避免极端观测直接以 $1/x_{t,i}$ 放大后破坏 FTRL 稳定性；clipping 保留极端观测达到阈值这一信息，用来更新尺度 $S_t$。若尺度也使用 skipped 损失，大损失被置零后阈值可能长期不增长；若累计损失直接使用 clipped 值，又会改变论文需要控制的估计误差结构。

\textbf{疑问二：skipping 使估计产生偏差，为什么算法仍可能取得遗憾保证？}

\textbf{解答：}分析并不假装偏差消失，而是把遗憾分成 FTRL 主项、尺度变化项和 skipping 误差。有限 $\alpha$ 阶矩给出尾部一阶矩界 $\sigma^\alpha C^{1-\alpha}$，阈值随尺度增长后，跳过部分的期望代价会下降。随机环境中，最优臂的截断非负性进一步使最优臂尾部项具有有利符号，从而只需控制非最优臂的尾部代价。

\textbf{疑问三：log-barrier 保证每个臂的概率为正，为什么还需要研究概率的乘法稳定性？}

\textbf{解答：}“概率为正”只排除了恰好为零，并不能阻止概率在相邻两轮间从很小变得更小，或在重要性加权后出现巨大变化。乘法稳定性比较加入一次估计前后的两个 FTRL 解，控制 $z_{t,i}/x_{t,i}$ 的比例；只有这种局部控制，才能把 Bregman 散度压到与 $(1-x_{t,A_t})^2$ 和处理后损失平方相关的量级。

\textbf{疑问四：论文为什么能用同一个算法同时得到对抗环境和随机环境的遗憾界？}

\textbf{解答：}对抗界主要依靠 FTRL 稳定性、有限矩尾界和自适应尺度，不使用固定 gap。随机环境额外提供 self-bounding 关系：探索非最优臂的概率总量可以由遗憾除以 $\Delta_{\min}$ 控制。证明在尺度足够大后把若干误差项吸收到遗憾自身，因此从同一套更新中得到 gap-dependent 对数界。环境适应发生在分析中，而不是算法先判断环境类型再切换策略。

\textbf{疑问五：“parameter-free”究竟表示什么，为什么证明里仍然会出现 $\alpha$、$\sigma$ 和 $\Delta_{\min}$？}

\textbf{解答：}它表示算法运行时不需要把未知的重尾阶数和矩尺度作为输入。性能上界当然仍由真实环境难度决定，所以这些参数会出现在定理和证明选择的阈值 $M$ 中。证明者用未知参数划分情况，不等于学习者在运行算法时知道这些参数；同时，parameter-free 也不表示算法不需要 $K$、时间范围或初始尺度等任何设置。

\subsection{8.15 下一步精读}

\textbf{建议独立完成的核验}
\begin{itemize}
\item 不看笔记推导命题 8.2、8.4、8.5 和 8.7，并独立手算一轮包含极端损失的 uniINF 更新。
\item 对照原文第 5 页逐行解释 Algorithm 1，区分输入、可观察量、分析中才定义的量。
\item 精读附录 C.1--C.3：证明概率乘法稳定性，检查隐式求导和 $(1-x)^2$ 的出现过程。
\item 精读附录 E 与 F：补齐 skipping 误差的停止阈值论证，以及完整 FTRL 分解。
\end{itemize}

\textbf{可进一步思考的问题（问题，不是已有成果）。}最优臂的截断非负条件在哪些数据中合理？若只做均值平移，尾部条件会怎样变化？能否改善对 $\Delta_{\min}$ 的依赖？在模拟实验中，如何同时记录 skip 比例、尺度 $S_t$、最优臂概率与遗憾，检查理论机制是否出现？这些问题可以作为后续与导师讨论的起点，但不应预先写成已完成的研究贡献。


\begin{center}\rule{0.5\linewidth}{0.5pt}\end{center}

\section{第九章
可重复实验规范}\label{ux7b2cux4e5dux7ae0-ux53efux91cdux590dux5b9eux9a8cux89c4ux8303}

理论学习应当与实验互相验证。一个基本实验至少记录：

\[
\texttt{seed, t, action, reward, instant\_pseudo\_regret}.
\]

\subsection{9.1 随机种子}\label{ux968fux673aux79cdux5b50}

固定 seed
并不是让每次调用随机函数都返回同一个数，而是固定整条伪随机序列。相同程序、参数和
seed 应产生相同结果。

环境 seed 与策略 seed 应分开：

\begin{itemize}
\tightlist
\item
  环境 seed 控制奖励序列；
\item
  策略 seed 控制随机探索和选臂。
\end{itemize}

这样才能区分``环境碰巧不同''和``策略自身随机性不同''。

\subsection{9.2
最低验收标准}\label{ux6700ux4f4eux9a8cux6536ux6807ux51c6}

\begin{enumerate}
\def\labelenumi{\arabic{enumi}.}
\tightlist
\item
  一条命令能够从头运行并生成原始CSV；
\item
  相同seed重复运行结果一致；
\item
  参数、时间范围和seed都能追溯；
\item
  累积pseudo-regret非负；
\item
  不只保存最终图片，也保存原始记录；
\item
  比较算法时使用相同环境seed集合。
\end{enumerate}

\subsection{9.3
推荐实验顺序}\label{ux63a8ux8350ux5b9eux9a8cux987aux5e8f}

\begin{enumerate}
\def\labelenumi{\arabic{enumi}.}
\tightlist
\item
  Gaussian环境比较Random、Greedy和\(\epsilon\)-greedy；
\item
  比较ETC与UCB在不同gap下的表现；
\item
  构造对抗奖励序列比较EXP3与随机策略；
\item
  构造重尾奖励，比较普通均值、截断均值与Median-of-Means；
\item
  将稳健估计器放入UCB，观察不同 \(\epsilon\) 下遗憾曲线的变化。
\end{enumerate}

\begin{center}\rule{0.5\linewidth}{0.5pt}\end{center}

\end{document}
